%!TEX TS-program = xelatex
\documentclass[]{../../../../friggeri-cv}

% High Assurance AI Programme Technical Lead variant; recruiter brief, 14 October 2026.
% Layout reference: Permutable AI variant; master and website PDF remain unchanged.

% Packages
\usepackage{pifont}
\usepackage{graphicx}
\usepackage{hyperref}
\usepackage{marvosym}
\usepackage{xcolor}
\usepackage{enumitem}
\usepackage{fancyhdr}

\graphicspath{ {../../../../images/} }
\definecolor{linkaccent}{HTML}{1E6291}
\definecolor{profileweb}{HTML}{B45309}
\definecolor{profilelinkedin}{HTML}{0A66C2}
\definecolor{profilegithub}{HTML}{24292F}
\definecolor{profilescholar}{HTML}{4285F4}
\newcommand{\profilelink}[3]{{\hypersetup{urlcolor=#1}\addfontfeature{Color=#1}\href{#2}{#3}}}
\hypersetup{
    colorlinks=true,
    urlcolor=linkaccent,
    linkcolor=linkaccent,
    citecolor=linkaccent,
    pdfborder={0 0 0},
    pdfauthor={Christos Hadjinikolis},
    pdftitle={Christos Hadjinikolis CV - High Assurance AI Programme Technical Lead}
}

% Canonical header override: restores a compact charcoal band while
% keeping a real space between first and last name for PDF text extraction.
\renewcommand{\header}[5]{%
  \begin{tikzpicture}[remember picture,overlay]
    \node [rectangle, fill=darkgray, anchor=north, minimum width=\paperwidth, minimum height=3.75cm] (box) at (current page.north){};
    \node [anchor=north] (name) at ([yshift=-0.72cm] box.north) {%
      \fontsize{36pt}{50pt}\color{white}%
      {\thinfont #1\ }{\bodyfont #2}%
    };
    \node [anchor=north] at ([yshift=-0.05cm] name.south) {%
      \fontsize{13pt}{20pt}\color{white}%
      \thinfont #3%
    };
    \node [anchor=north] at ([yshift=-0.55cm] name.south) {%
      \fontsize{11pt}{18pt}\color{white}%
      \hypersetup{urlcolor=white,linkcolor=white}\thinfont #4%
    };
    \node [anchor=north] at ([yshift=-3.08cm] box.north) {%
      \fontsize{10.5pt}{16pt}\color{white}%
      \hypersetup{urlcolor=white,linkcolor=white}\thinfont {\Letter\hspace{0.2em} #5}%
    };
  \end{tikzpicture}
  \vspace{2.2cm}
  \vspace{-2\parskip}
}

\renewcommand{\entry}[4]{%
  \parbox[t]{\dimexpr\textwidth-\entryleftinset\relax}{%
    \textbf{#2}%
    {\footnotesize\addfontfeature{Color=lightgray}\hspace{0.6em}#1}%
    \hfill%
    {\small\color{textcolor} #3}\\%
    #4\vspace{\parsep}%
  }\\}

\renewcommand{\entrytwo}[4]{%
  #1&\parbox[t]{14cm}{%
    \textbf{#2}%
    \hfill%
    {\small\color{textcolor} #3}\\%
    #4\vspace{\parsep}%
  }\\}

% Set up footer
\pagestyle{fancy}
\fancyhf{}  % Clear header/footer
\fancyfoot[C]{\small Page \thepage}
\renewcommand{\headrulewidth}{0pt}
\renewcommand{\footrulewidth}{0pt}

% XeLaTeX image fix
\ifxetex
  \usepackage{letltxmacro}
  \setlength{\XeTeXLinkMargin}{1pt}
  \LetLtxMacro\SavedIncludeGraphics\includegraphics
  \def\includegraphics#1#{%
    \IncludeGraphicsAux{#1}%
  }%
  \newcommand*{\IncludeGraphicsAux}[2]{%
    \XeTeXLinkBox{%
      \SavedIncludeGraphics#1{#2}%
    }%
  }%
\fi

\begin{document}
\header
    {Christos}
    {Hadjinikolis}
    {Engineering Manager | Reliable AI Systems \& Technical Leadership}
    {Ph.D. Computer Science | UCL Associate Researcher | AI Standards}
    {\small \href{mailto:christos.hadjinikolis@gmail.com}{christos.hadjinikolis@gmail.com}}

\begin{aside}
    \section{profiles}
        {\footnotesize \profilelink{profileweb}{https://christos-hadjinikolis.github.io/}{website}\\\profilelink{profilelinkedin}{https://uk.linkedin.com/in/christoshadjinikolis}{LinkedIn}\\\profilelink{profilegithub}{https://github.com/Christos-Hadjinikolis}{GitHub}\\\profilelink{profilescholar}{https://scholar.google.com/citations?hl=en&user=Rr3KTx8AAAAJ}{Google Scholar}}
    \section{leadership}
        {\footnotesize 6 direct reports\\10-person team\\hands-on delivery\\hiring/coaching}
    \section{systems}
        {\footnotesize real-time data\\Kafka/Flink\\model/API serving\\observability}
    \section{ML assurance}
        {\footnotesize model/data versioning\\evaluation/replay\\data contracts\\accuracy monitoring}
    \section{maturity}
        {\footnotesize ADRs/forums\\runbooks/on-call\\local E2E tests\\rollback/fallback}
    \section{domain}
        {\footnotesize AI standardisation\\auditability\\symbolic reasoning\\explainability}
    \section{tooling}
        {\footnotesize \textbf{PyPI}: \href{https://pypi.org/project/dynamicio/}{dynamicio}\\\href{https://pypi.org/project/skeleton-replay/}{skeleton-replay}\\\href{https://plugins.jetbrains.com/plugin/32807-skeleton-replay}{JetBrains plugin}}
\end{aside}

\section{about me}
{\normalsize I am a \textbf{hands-on Engineering Manager and ML systems lead}, combining production AI engineering with doctoral research and AI standards work. I lead teams turning research into dependable services, with model/data versioning, evaluation, replay and monitoring built into delivery. Alongside managing people, I continue substantial individual-contributor work on production code and ML/data pipelines. My UCL research and teaching and AI committee work focus on auditability, explainability and responsible adoption.}
\section{work-experience}
\begin{entrylist}
    \entry
    {12/2020--present}
    {Vortexa Ltd, London, UK}
    {Engineering Manager | ML Systems Lead}
    {
        {\small
        \setlength{\parindent}{0pt}
        \setlength{\leftskip}{0pt}
        I lead technical direction and delivery for Vortexa's live post-ingestion ML/data estate, turning \textbf{~6M} vessel-position records/hour into production commodity intelligence for \textbf{13.5K} monitored vessels. Manage \textbf{6} direct reports and lead a \textbf{10-person} cross-functional team accountable for model quality, data-product interfaces, client and stakeholder clarity, reliability, and delivery maturity. Joined as Senior ML Engineer in December 2020; Engineering Manager / ML Systems Lead since February 2022.

        }
        \begin{itemize}[leftmargin=1.1em]
            \item \footnotesize\textbf{People \& Delivery:} Manage \textbf{6} direct reports across Product, SME analysis, Data Science, and Data Engineering; own performance reviews, career development and delivery. Hiring manager for \textbf{6} roles; shaped onboarding and hiring practices as Data Production grew from \textbf{4 to 30+}. Retained the full team and supported every member's promotion or progression (July 2026 snapshot).
            \item \footnotesize\textbf{Technical Direction:} Shape architecture across model services, data pipelines and evaluation/replay interfaces. Resolve model-quality questions with Product, researchers and domain experts, translating competing priorities into measurable objectives and delivery plans.
            \item \footnotesize\textbf{Production ML Systems:} Led research-to-production delivery for destination and arrival-time sequence/transformer models in \textbf{PyTorch}; established MLflow, model/data versioning, evaluation, replay, monitoring and model-serving workflows.
            \item \footnotesize\textbf{Hands-on Implementation:} Led the \textbf{Kafka Streams-to-Flink} migration: researched trade-offs, wrote the proposal and personally built the first replacement application into production. Improved operational visibility and compute scaling independent of Kafka partitioning.
            \item \footnotesize\textbf{Operational Assurance:} Established live accuracy monitoring, signal-quality controls, runbooks, rollback/fallback paths and shared on-call practices, so teams could investigate prediction failures and operate evolving models reliably.
            \item \footnotesize\textbf{Engineering Maturity:} Standardised I/O seams, schema/data validations, repo archetypes, ADRs, dev containers, local E2E tests, and DORA/Jira delivery-health signals; cut a \textbf{3h} DAG feedback loop to under \textbf{5} min.
            \item \footnotesize\textbf{Shared Engineering Practice:} Started a weekly architecture forum to share decisions and reduce knowledge silos. Ran stakeholder workshops and design reviews to make trade-offs explicit and align engineering work with product and client needs.
        \end{itemize}
    }
    \entry
    {04/2016--12/2020}
    {DataReply, London, UK}
    {Senior Consultant | Data Scientist / ML Engineer}
    {
        {\small
        \setlength{\parindent}{0pt}
        \setlength{\leftskip}{0pt}
        Joined Data Reply's London spin-off as its first consultant, operating inside and contributing to a practice that grew from \textbf{three practitioners and two managers to 30+ consultants}. Progressed to Senior Consultant; shaped opportunities, agreed deliverables/timelines, managed expectations, placed consultants, and helped build the consulting culture around clear technical delivery across UBS, Vodafone, CNHi, and other clients.

        }
        \begin{itemize}[leftmargin=1.1em]
            \item \footnotesize\textbf{Vodafone Group:} Led multi-workstream GCP data-science platform work based on Kubeflow; drove requirements, migration planning, CI/CD, workflow automation, ML lifecycle support, and feature-engineering migration.
            \item \footnotesize\textbf{CNHi:} Lead Data Scientist/Scrum Master on predictive maintenance (Azure/Databricks); translated stakeholder needs into DS/DE backlog and guided a \textbf{PySpark} team over live telemetry.
            \item \footnotesize\textbf{UBS:} Embedded consultant who grew from Data Scientist into ML Engineer through XP/pairing while delivering graph/process-mining analytics, observability, and real-time \textbf{Kafka}/Elasticsearch pipelines.
        \end{itemize}
    }
    \entry
    {2010--2016}
    {KCL, UCL, David Game College, UK}
    {Teaching and Technical Training}
    {
        {\small
        \setlength{\parindent}{0pt}
        \setlength{\leftskip}{0pt}
        Delivered courses and tutorials in algorithm design, data structures, databases/SQL, and Java/Python during and after doctoral research; developed presentation discipline, learner support, and mentoring habits.

        }
    }
\end{entrylist}
\newpage
\setlength{\parskip}{2pt}
\setlength{\parsep}{0pt}
\specialsection{education}
\begin{entrylist2}
    \entrytwo
        {\small 2010--2014}
        {\small Ph.D. in Computer Science}
        {\small King's College London, UK}
        {\small \emph{\textbf{Thesis}: Persuasion Dialogues \& Opponent Modelling.}

        {\footnotesize Developed logical inference methods for large \textbf{Knowledge Graphs}, combining symbolic reasoning, \textbf{graph analytics}, and Bayesian techniques (e.g., MCMC). Received \textbf{Best Poster Award} at IJCAI 2013 (413 submissions), \textbf{EPSRC PhD Scholarship}, \textbf{Outstanding TA Award}, and \textbf{Graduate Certificate in Academic Practice}.
        \\ \textbf{\textcolor{blue}{\href{https://scholar.google.com/citations?hl=en&user=Rr3KTx8AAAAJ}{\underline{Access Publications on Google Scholar}}}}}
        }
    \entrytwo
        {\small  2004--2010}
        {\small Diploma (BEng) in Computer Engineering}
        {\small University of Thessaly, GR}
        {\footnotesize Five-year polytechnic degree in Computer Science and Engineering with a \textbf{strong focus on mathematics} (10 courses), covering theoretical aspects, networking, and telecommunications. \emph{\textbf{Majored in Artificial Intelligence}}.}
    \entrytwo
        {\small  2002--2004}
        {\small Military Officer: Second Lieutenant}
        {\small National Guard, CY}
        {\small Completed national service as a commissioned Second Lieutenant.}
\end{entrylist2}
\specialsection{research \& AI standards}
\begin{entrylist2}
    \entrytwo
        {\footnotesize since 10/2024}
        {\small University College London, Dept. of Information Studies}
        {\small Associate Researcher}
        {\small Research/teach practical AI standardisation: AI Act, auditability, model/data versioning, explainability, and safe adoption.}
    \entrytwo
        {\footnotesize since 2021}
        {\small ISO - International Organisation for Standardisation}
        {\small Committee Expert Member for JTC'21}
        {\footnotesize Member of JTC'21 Working Group 3, developing AI standards aligned with EU policy and international norms.}
\end{entrylist2}
\specialsection{selected AI/platform tooling}
\begin{entrylist2}
    \entrytwo
        {\small}
        {\small \href{https://christos-hadjinikolis.github.io/2026/08/01/skeleton-replay-runtime-architecture-evidence.html}{Promet}}
        {\small Personal project | Local AI harness}
        {\footnotesize Personal AI assistant runtime exploring controlled tool execution, schema validation, approval boundaries, durable state, traces and replay/evaluation. Hands-on work on the application controls around model behaviour, with explicit permissions and inspectable execution evidence.}
    \entrytwo
        {\small PyPI}
        {\small \href{https://pypi.org/project/skeleton-replay/}{skeleton-replay}}
        {\small Runtime evidence}
        {\footnotesize Python developer tool that turns script/pytest runs into traces, architecture snapshots, workflow evidence, and replayable reports for review, debugging, onboarding, and evidence-led code understanding.}
    \entrytwo
        {\small JetBrains}
        {\small \href{https://plugins.jetbrains.com/plugin/32807-skeleton-replay}{Skeleton Replay plugin}}
        {\small IDE workflow}
        {\footnotesize Free PyCharm/IntelliJ plugin bringing Skeleton reports, workflow evidence, and source navigation into the IDE; reflects a bias for developer experience, evidence, and shared understanding.}
    \entrytwo
        {\small PyPI}
        {\small \href{https://pypi.org/project/dynamicio/}{dynamicio}}
        {\small Data contracts}
        {\footnotesize Created a Python library for explicit I/O boundaries, schema/data validation and local/dev/prod dataset switching. Used across 15+ repositories to make ML/data workflows easier to test and review.}
\end{entrylist2}
\specialsection{talks/interviews}
\begin{entrylist2}
    \entrytwo
        {\small 2023}
        {\small Bill Raymonds, Agile in Action}
        {\small Podcast Interview}
        {\small \textbf{Agile Data Science}: Sailing Through Data Science: The Agile Journey at Vortexa, \underline{\textbf{Details:} \href{https://christos-hadjinikolis.github.io/2023/10/31/Agile-In-Action.html}{YouTube}}}
    \entrytwo
        {\small 2022}
        {\small Open Data Science Conference}
        {\small Industry Conference Talk}
        {\small \textbf{dynamicio}, a published \href{https://pypi.org/project/dynamicio/}{\textbf{PyPI library}} for abstracting I/O in ML/data workflows, \underline{\textbf{Details:} \href{https://odsc.com/speakers/dynamicio-a-pandas-i-o-wrapper-why-you-should-start-your-ml-ops-journey-with-wrapping-your-i-o/}{here}}}
    \entrytwo
        {\small 2020}
        {\small iunera: Big Data Innovators}
        {\small Interview Blog Post}
        {\small The Agile approach in data science explained by an ML expert, \underline{\textbf{Details:} \href{https://www.iunera.com/kraken/interviews/the-agile-approach-in-data-science-explained-by-an-ml-expert/}{here}}}
    \entrytwo
        {\small 2020}
        {\small Big Data Warsaw}
        {\small Industry Conference Talk}
        {\small Monitoring \& Analysing Communication and Trade Events as Graphs, \underline{\textbf{Details:} \href{https://bigdatatechwarsaw.eu/agenda/}{here}}}
    \entrytwo
        {\small 2018}
        {\small Connected Data London}
        {\small Panelist @ Industry Conference}
        {\small AI tribes and interconnections: what's graph got to do with it?, \underline{\textbf{Details:} \href{http://connected-data.london/2018/09/23/ai-tribes-and-interconnections-whats-graph-got-to-do-with-it/}{here}}}
    \entrytwo
        {\small 2018}
        {\small Minds Mastering Machines}
        {\small Industry Conference Talk}
        {\small A talk on doing Data Science the Agile Way, \underline{\textbf{Details:} \href{https://2018.mcubed.london/sessions/data-science-agile-way/}{here}}}
\end{entrylist2}
\specialsection{certificates \& awards}
\begin{entrylist2}
    \entrytwo
        {\small 2017--2020}
        {\small Selected ML/cloud credentials}
        {\small AWS, Google Cloud, O'Reilly}
        {\footnotesize AWS ML Specialty; Google Professional Data Engineer; Apache Spark Developer.}
    \entrytwo
        {\small 2016--2019}
        {\small Data/graph engineering credentials}
        {\small Coursera, Elastic, Neo4j}
        {\footnotesize Process Mining; Graph Analytics; Advanced Elasticsearch; Neo4j Certified Professional.}
\end{entrylist2}
\end{document}
